\section{通过项定义的偏函数}

\begin{metatheorem}{通过项定义函数}{}
	设$T,A$是项,$x,y$是不同的字母,其中$x$不在$A$中出现,$y$不在$T$或$A$中出现,$R$是公式$x\in A\wedge y=T$.那么
	\begin{enumerate}
		\item $R$有唯一一个相对$x$和$y$的二元关系$F$,并且$F$是偏函数关系.
		\item $\pr_{1}\left(F\right)=A,\pr_{2}\left(F\right)=\left\{T\middle|x\in A\right\}$.
		\item 对任一$x\in A$有$F\left(x\right)=T$.
	\end{enumerate}
\end{metatheorem}

\begin{remark}{记号说明}{}
	记号如上.设$C$是包含$\left\{T\middle|x\in A\right\}$的集合.记号没有歧义时,上述元定理确定的偏函数$\left(F,A,C\right)$可记作下述形式之一:
	\begin{gather*}
		x\mapsto T\left(x\in A\right)\\
		\left(T\left(x\right)\right)_{x\in A}\\
		x\mapsto T
	\end{gather*}
	我们甚至可以直接用$T$来代表该函数.
\end{remark}

\begin{example}{映射的表达式展开}{}
	设$f:A\to B$是映射,则$f$和映射$f:A\to B,x\mapsto f\left(x\right)$相等,后者有时候简单写作$x\mapsto f\left(x\right)$或$\left(f\left(x\right)\right)_{x\in A},\left(f_{x}\right)_{x\in A}$.
\end{example}

\begin{example}{投影映射}{}
	设$G$是二元关系.定义在$G$上的映射
	\begin{gather*}
		G\to\pr_{1}\left(G\right),z\mapsto\pr_{1}\left(z\right)\\
		G\to\pr_{2}\left(G\right),z\mapsto\pr_{2}\left(z\right)
	\end{gather*}
	在不引起混淆的情况下,分别记作$\pr_{1}$和$\pr_{2}$,分别称为第一个投影和第二个投影.
\end{example}



























